\documentclass{article}
\usepackage{enumerate}
\usepackage[margin=1.4in, top=1in]{geometry}
\title{Teaching Statement}
\author{Joel Villatoro}
\date{}

\usepackage{blindtext}
\usepackage{hyperref}
\begin{document}
\vspace{-0.5in}
\maketitle

\section{Introduction}

If there is anything the past two years as a lecturer at Indiana University have taught me its that I like teaching.
I've always thought that enjoying a job is a prerequisite to doing that job well.
This is true enough for me at least.

There are a few specific things I like about it.
For one I've always enjoyed the art of finding a good explanation. 
Finding just the right way to communicate an idea to another person is a joy. 
It can be quite a puzzle and, as a mathematician, I quite enjoy solving such problems.

The other thing I like about it is the sense of connection. When a person learns from another person, there is a natural bond that forms. 
A bond that is added to a continuous chain of knowledge that humans have been passing down since time immemorial. 
I think this sense of shared knowledge and understanding between people is incredibly valuable.
Not just because it can help us solve difficult problems in the future, but because the the connections we form are ultimately what hold our society together.

OK, hopefully that is enough poetry. Ultimately, teaching is not just an art but also a craft. Let me know explain how I think about and approach my craft.

\section{Classroom Practice}

The classroom is the centerpiece for the educational experience of the student. It's not everything, but it forms the core of each student's experience with a class.
I approach in-class instruction with the assumption that the class will be the student's first exposure to the material. 
The goal of the teacher is to facilitate a student experience that maximizes the impression that material will make.
Although I am not a researcher of education per se, my understanding is that the best evidence show that ``interactive'' and ``active learning'' approaches to class time are the most effective.

My default approach is to provide students with prepared note "blanks" in advance which are to be filled by the students during the lecture. 
The note blanks are designed ot be a hybrid between traditional lecture notes and a worksheet.
The typical procedure is a mix of the following:
\begin{itemize}
\item I work out examples while students follow along. I try to verbalize my thought process in detail as the problem gets worked out. 
I also try to call specific attention to how solutions are organized and encourage students to model them closely.
\item Students work out examples. Time is set aside for students to attempt the problem on their own. If the physical classroom format permits it, students may be grouped and encouraged to collaborate in groups.
\item Additional activities to break the routine. The details of how this works varies significantly from course to course. However, some examples are:
\begin{itemize}
\item In the class ``Mathematics of Decision and Beauty" (100-level class for majors with minimal math requirements): At various times throughout the semester I created small websites that had interactive demos that allowed students to interact with course material a variety of ways.
\item In the class ``Brief Survey of Calculus" (100-level class, comparable to ``Business Calculus" at many institutions.): Various activities such as: Desmos based demos, listening to a song that references a mathematical formula, group quizzes.
\end{itemize}
\end{itemize}
\subsection*{Student Feedback on Classroom Practices}
\begin{itemize}
\item "[He] would start each class by explaining the concepts, and then would spend the rest of class going through examples, which was a really good way to learn Calculus. He would often challenge us to help him solve new questions so that would keep us engaged." (M211)
\item "[He] always has more than one way of explaining material to students. So, if you're confused he switches over to a different method to help you. I also like the course packets every week because it was easy to follow along with class and write down what the professor was teaching in real time." (M119)
\end{itemize}



% As an academic mathematician studying pure mathematics, I consider teaching and mentoring as the primary way I can leave a lasting positive impact. 
% I also consider it a deeply rewarding part of my job.
% The core of my teaching philosophy can be described by my two main focuses, student engagement and accessibility. 
% By student engagement, I mean getting students involved in the course beyond a lecture/homework format. 
% Accessibility means remaining available and approachable to my students and mentees.

% My experiences teaching cover a wide range of ages, skill levels, backgrounds, and formats. 
% However, I believe that I still have much to learn and great teaching and mentorship is not a specific goal but a lifelong pursuit. 
% At the University of Illinois I was successful as a TA and was awarded for excellence. 
% At KU Leuven in Belgium I was faced with a significantly different student body and classroom environment and that presented me with many new challenges. 
% At Washington University, due to restrictions placed by my NSF grant, I was limited in the amount of paid teaching I could do but I did teach a large lecture course as well as an upper division course with a smaller class size. In addition to classroom work, at all of these institutions I have also engaged with undergraduates in a mentor/advising context. I will first discuss my views and philosophy regarding classroom instruction and then I will discuss my experiences as a mentor.

% My views on how to maintain an engaging classroom have been shaped by my experiences. 
% My formative teaching experience was as a teaching assistant at the University of Illinois. 
% At UIUC, the main approach was active learning, which meant that the students were split into small groups and given a worksheet. 
% My role was as a facilitator for discussions and provide students with hints and suggestions if they get stuck. In this time I realized the importance of learning to identify struggling or distracted students and to take proactive action rather than simply waiting for questions. In my experience, this format was generally very successful and, at first, it was the only teaching format I was familiar with.

% Unsurprisingly, when it came time for me to work as the primary instructor, I found that maintaining student engagement was significantly more challenging. 
% Thus far, I have employed a few strategies to break up the tedium of lectures and improve student engagement. 
% My favorite and generally most successful method is the creation of interactive visualizations of mathematical concepts. 
% For example, while teaching differential equations, I created an interactive phase diagram visualizer\footnote{link: \url{https://www.desmos.com/calculator/tj3vn0rcqj}]} so that students could develop intuition for integral curves and eigenvectors. Another example is python notebook that illustrates a topological proof of the fundamental theorem of algebra \footnote{link:\url{https://colab.research.google.com/drive/1A898TEOGSHCNcYKM0BWP8YmsE2PLFz9M?usp=sharing}} which I created for another instructor. 
% I am also frequently on the lookout for online resources (such as videos or other multimedia) that can compliment the lectures and worksheets. 
% Especially for undergraduate level topics, there is a vast quantity of high quality visualizations and instructional resources. 
% I intend to continue to experiment with and learn new techniques to provide an engaging student-centered classroom.

% In any mathematics course, the majority of learning will inevitably take place by doing exercises and peer discussion. 
% For courses with large enrollment, this means that the exercise sections and teaching assistants have a particularly significant impact on student outcomes. 
% When I taught differential equations the students reported that they felt they learned the most from doing the worksheets I wrote and from interacting with their TA. 
% While this is not surprising, it highlights the fact that primary instructors for such classes have their work cut out for them when it comes to directly impacting student success. 
% Many of the levers they have available to them are fundamentally indirect such as preparing course materials or supervising and training TAs.

% When teaching a class, I try to frequently remind students that they are welcome to come to my office and email me with their concerns.
% However, I have found that this is not always enough as some students require more proactive measures to ensure that they are getting the help they need. 
% After each exam, I try to identify which students appear to be struggling and I send them a personal invitation to visit me to discuss the class. 
% It's important to approach these interactions with sensitivity because students can often feel embarrassed or demoralized when they don't achieve the success they envisioned for themselves. 
% For example, I avoid using phrases such as "you appear to be struggling in the class" as it can be seen as overly personal and assigning blame to the student. 
% I also try to remain cognizant of and sensitive to the diverse backgrounds of my students.
% This means taking advantage of common ground I may have with other students as well as showing respect for their cultural and personal experiences of my students.

% While there are overlapping skills between mentorship and classroom instruction. 
% I would like to now take the time to discuss some of my views on and experiences in mentoring/advising students. 
% Broadly, my philosophical takeaways from these experiences is that, when mentoring any student, it is important to take into account and focus on the student's goals for themselves. 
% I have also found that, when working with undergraduates, one of the best ways to connect with my mentees is to learn alongside them.

% At UIUC, I worked a mentor for undergraduates through the Illinois Geometry Lab. 
% The Illinois Geometry Lab is a program run by the mathematics department which pairs small groups of undergraduates with a faculty member and graduate student. 
% The faculty member provides the undergraduates with a project, usually a visualization of some mathematical concept, while the graduate student provides regular support while the undergraduates work on the project. 
% I worked on multiple projects during my time but the first one involved numerical simulations of differential equations. At the time, I had little background in programming or computational mathematics. 
% Despite this, I found that this was in many ways a benefit as I was better able to bond and connect with the undergraduate students by learning along side them. 

% At KU Leuven my primary work as a mentor was supervising a Bachelor project. 
% This involved advising a group of undergraduate mathematics majors while they studied chosen topic in mathematics. 
% This experience was distinct from the IGL as the students were mathematics majors primarily interested in pursuing a career in pure mathematics. 
% Much of our activities were reading and discussing mathematics articles and excerpts from textbooks. 
% My primary role was to guide and advise them as they dipped their toe into reading and learning about topics in more advanced mathematics. 
% The specific topic we investigated was crystallography and Bieberbach groups which was not a topic I was directly familiar with. 
% Once again, I was reminded of how rewarding it can be to learn alongside my students.

% Overall, I feel that my past mentorship experiences have provided me with the tools to be able to flexibly work with undergraduates and graduate students in a way that takes into account their personal interests and objectives. 
% I hope to be able  to continue this work at your university.
\end{document}
